\begin{problem}[第34届物理竞赛复赛][回旋加速器]{76}
    某种回旋加速器的设计方案如俯视图 a 所示，图中粗黑线段为两个正对的极板，其间仅在带电粒子经过的过程中存在匀强电场，两极板间电势差为 $U$。两个极板的板面中部各有一狭缝（沿 $OP$ 方向的狭长区域），带电粒子可通过狭缝穿越极板（见图 b）；两细虚线间（除开两极板之间的区域）既无电场也无磁场；其它部分存在匀强磁场，磁感应强度方向垂直于纸面。在离子源 $S$ 中产生的质量为 $m$、带电量为 $q$（$q > 0$）的离子，由静止开始被电场加速，经狭缝中的 $O$ 点进入磁场区域，$O$ 点到极板右端的距离为 $D$，到出射孔 $P$ 的距离为 $bD$（常数 $b$ 为大于 $2$ 的自然数）。已知磁感应强度大小在零到 $B_{\max}$ 之间可调，离子从离子源上方的 $O$ 点射入磁场区域，最终只能从出射孔 $P$ 射出。假设如果离子打到器壁或离子源外壁便即被吸收。忽略相对论效应。求
    \begin{subproblem}
        可能的磁感应强度 $B$ 的最小值；
    \end{subproblem}
    \begin{subproblem}
        磁感应强度 $B$ 的其它所有可能值；
    \end{subproblem}
    \begin{subproblem}
        出射离子能量的最大值。
    \end{subproblem}
\end{problem}